% =====================================================================
% PAPER 8: Temporal Stability of LLM Personas (Gonnermann-Müller)
% PERSONA_CATEGORY: Surrogate user persona (LLM-simulated participant)
% PERSONA_SUBCATEGORY: Persona-as-simulated-research-subject (social-science simulation agent)
% PERSONA_SUBCATEGORY: Persona-as-psychometrically-measurable trait configuration (stability/drift construct)
% =====================================================================

\subsection*{Paper 8 --- \textit{Temporal Stability of LLM Personas: A Prerequisite for Social Science Simulation} (Gonnermann-Müller)}

This paper uses \emph{persona} in the computational social-science sense: a persona is a prompt-encoded specification of ``characteristics, behavioral tendencies, and goals'' that configures an LLM agent so that it can stand in for a human subject in simulated studies. The persona is therefore neither a design artifact nor a communicative character for end users; it is the \emph{experimental unit of a synthetic participant}, the thing whose fidelity determines whether simulated trust dynamics, opinion change or group interaction count as evidence about people at all. Notably, the persona here is not a rich narrative sketch of a named individual but a graded, clinically anchored trait configuration: personas are operationalized as four ADHD intensity conditions (high, moderate, low, and a no-persona default baseline), realized through three semantically equivalent prompt variants (text-based and scale-based formulations) across seven proprietary and open-source LLMs.

The paper's distinctive contribution is to treat persona as something that must be \emph{measured}, and specifically measured \emph{over time}. Borrowing the multi-informant paradigm from clinical psychology, the author separates the persona's \emph{self-reported} identity (the target model completing the 12-item ADHD Index of the Conners' Adult ADHD Rating Scales in first person) from its \emph{observer-rated behavioural expression} (three independent LLM judges scoring only the behavioural evidence in generated workday narratives and conversation transcripts, without access to the persona instructions). Persona stability is then decomposed into two temporal dimensions: between-conversation stability (50 independent instantiations per condition) and within-conversation stability (18-turn dialogues probed at turns 6, 12 and 18). The results introduce a conceptual dissociation that matters for how the field defines persona: self-reported persona identity remains stable, while enacted persona expression decays (high-intensity personas drop $\sim$20\%, $M=17.5 \rightarrow 14.0$, with rising variance $SD=2.49 \rightarrow 3.87$). In other words, the paper implicitly distinguishes \emph{persona-as-claimed-identity} from \emph{persona-as-observable-behaviour}, and warns that only the latter carries validity for interaction simulation.

Persona is thus framed primarily as a \emph{validity risk object} and a governance target: drift toward an ``average persona'' (attributed to context-window dilution and RLHF pressure toward normative, socially desirable output) would make emergent patterns in multi-agent simulations artifacts of model inconsistency rather than social phenomena. The paper's responsible-AI framing is explicit --- personas representing accessibility needs or novice users that silently lose those characteristics mid-session produce ``unreliable or potentially harmful'' insights --- and it proposes community standards (repeated mid-session consistency checks, reporting of variability metrics, replication across seeds and configurations, documentation of model/prompt/temperature choices). This paper stretches the existing ``surrogate user persona'' label away from usability testing toward the simulation of \emph{human subjects and clinical populations} generally; the shared commitment is persona-as-proxy-for-a-person, but the evaluative criterion here is psychometric reliability rather than design usefulness.

% ---------------------------------------------------------------
% Updated running taxonomy table (Papers 1--8)
% ---------------------------------------------------------------

\begin{table*}[t]
\centering
\caption{Running taxonomy of persona conceptualizations across workshop papers (updated through Paper 8).}
\label{tab:persona-taxonomy}
\small
\begin{tabular}{p{0.5cm} p{3.6cm} p{3.4cm} p{3.6cm} p{3.6cm}}
\toprule
\textbf{\#} & \textbf{Paper} & \textbf{Main category} & \textbf{Subcategory} & \textbf{Operationalization / unit} \\
\midrule
1 & AI Personas for UX Research in Large Scale Retail (Sharma et al., Walmart) &
Surrogate user persona (LLM-simulated participant) &
(a) Persona-as-test-participant; (b) Persona-as-latent-trait-vector (activation steering) &
Demographic/behavioural archetype profiles prompted into a multimodal LLM to produce synthetic usability feedback; validated against 150 human participants. Secondary: persona as steering direction in BLIP-2 activations \\
\addlinespace
2 & AI Personas in Highly Regulated Environments: Tools for Calibrating User Reliance? (Vasiliu \& Guarese) &
System-facing agent persona (deployed AI character) &
Persona-as-interaction-metaphor for trust/reliance calibration (incl.\ adaptive / metaphor-fluid personas) &
Two prompt-encoded conversational personas (\textit{Expert Operator} vs.\ \textit{Machine}) defining tone, stance and dialogue structure of a voice CAI; manipulated as an experimental variable and evaluated via trust/workload/usability measures and a behavioural disclosure probe \\
\addlinespace
3 & Evaluating Conversational Agent Integration in Peer Support Groups (Uetova et al.) &
Surrogate user persona (LLM-simulated participant) &
Persona-as-data-derived individual profile (behavioural clone) for multi-agent social simulation; explicitly contrasted with large ``persona catalogues'' (NLP persona datasets) &
One persona per real past participant: LLM-generated summary of that person's chat history, demographics and survey answers (tone, emoji use, thematic focus), used to condition synthetic peer replies in a replayed group chat; combined with probabilistic responder selection and ``buddy'' pairing; framed as a pre-deployment testbed for agent prompt/logic refinement, with stated limits on fidelity and inclusivity \\
\addlinespace
4 & Grounding Persona-Driven Usability Simulation in Established Standards (Touir, Barika Ktata \& Soui) &
Surrogate user persona (LLM-simulated participant) &
Persona-as-autonomous-task-agent (embodied evaluator in a perceive--decide--act loop); Persona-as-edge-case/accessibility probe &
Supervisor-Agent auto-synthesises personas from UI purpose + source code using a constrained schema (goal, tech-literacy level, accessibility constraint); each persona conditions a VLM/LLM User-Agent that clicks, scrolls and types on a live prototype, emitting action traces, error logs and think-aloud self-reports mapped to ISO 9241-110, Nielsen's heuristics and WCAG. Named exemplars (``Martha R.,'' ``David L.'') act as edge-case probes; governance layer forbids unsupported demographic stereotyping and requires provenance and human escalation \\
\addlinespace
5 & Ideal Persona AI in Real-time (Pham Thi Ngoc Anh) &
Self-expression persona (user's idealized self rendered by GenAI) &
(a) Persona-as-generated-self-image (real-time text-to-image self-portrait); (b) Persona-as-ethical-behaviour-envelope (transparent ``default ethical persona'' of the system) &
Persona defined as ``a personality belonging to a participant'' --- an ideal self co-authored with GenAI. Six designers wrote ideal-persona descriptions with ChatGPT, then fed five keywords/questions as prompts to StreamDiffusion in TouchDesigner and watched a live image of that persona; screen capture, field notes, think-aloud and post-study questions analysed thematically. Breakdowns: inconsistent persona prototypes, uncanny valley, disconnection/trust erosion. Proposes moving persona design from ``creative writing'' to ``structural engineering'': prompt training, declared default ethical persona, ethical constraints exposed as UI/UX controls, continuous user control in real time \\
\addlinespace
6 & Mapping Personas to Text Transformations: A Taxonomy Outline for Content Adaptation (Sillano, De Russis, Cal\`o, Troncy \& Lisena) &
Content-adaptation audience persona (NLP prompt-conditioning target reader) &
(a) Persona-as-decomposable-trait-vector mapped to auditable text-transformation operators; (b) Persona-as-target-reader-profile for content personalization &
Persona = specification of the intended reader that conditions an LLM rewrite, critiqued as a ``black-box prompt modifier''. Decomposed into four textually-correlated dimensions (demographic, domain knowledge, cognitive, contextual) synthesised from role-playing/personalization surveys; each dimension mapped in a taxonomy table to named operators (lexical, structural, content, style, tone) with worked examples. Two tabular persona profiles (P1 novice student, P2 science journalist) drive a five-step operator chain rewriting a thermodynamics sentence with GPT-5.4, each edit traceable to an operator + persona dimension. Goal: transparency, bias auditability and writer agency; mappings not yet empirically validated \\
\addlinespace
7 & \textsc{Synonymix}: Privacy-Preserving Life Story Personas via Graph-Based Abstraction (Chen \& Kumar) &
Surrogate user persona (LLM-simulated participant) &
(a) Persona-as-narrative-life-story profile (high-fidelity behavioural proxy for a real individual); (b) Persona-as-privacy-preserving synthetic composite (graph-sampled from many real people); (c) Persona-as-collective artifact (``unigraph'' for meso-level sensemaking) &
Persona framed as data grounding for a generative agent, on a fidelity continuum from ``flattened'' demographic personas to McAdams-style life story personas. Each persona is converted into a typed knowledge graph (Subject / Factual / Interpretive nodes, with I-nodes produced by LLM ``expert reflection''), genericised and semantically deduplicated into a merged \emph{unigraph}; new personas are sampled by a Thematic Random Walk anchored on an interpretive node. Evaluated as three 30-agent banks (Demographic, Life Story, Synonymix) answering 108 GSS 2022 items, scored by EMD/TVD distributional distance ($p<0.001$, $r=0.59$ on ordinal items) and by Maximum Source Contribution ($\overline{MSC}=0.129$) as a structural privacy metric; DP over node histograms deferred. Persona additionally reframed as a shareable collective representation supporting cross-temporal sensemaking and consensus building rather than prediction \\
\addlinespace
8 & Temporal Stability of LLM Personas: A Prerequisite for Social Science Simulation (Gonnermann-Müller) &
Surrogate user persona (LLM-simulated participant) --- here broadened to \emph{simulated human research subject} &
(a) Persona-as-simulated-research-subject (social-science simulation agent); (b) Persona-as-psychometrically-measurable trait configuration (stability/drift construct); (c) Implicit split between persona-as-claimed-identity (self-report) and persona-as-enacted-behaviour (observer rating) &
Persona = prompt-encoded set of characteristics, behavioural tendencies and goals configuring an LLM agent to stand in for a human subject. Operationalized as four graded conditions on a clinical dimension (high / moderate / low ADHD + no-persona default), instantiated via three semantically equivalent prompt variants across seven LLMs. Stability measured with a multi-informant design: first-person self-report on the 12-item Conners' ADHD Index vs.\ blind ratings by three LLM observers scoring behaviour only. Experiment I: 50 independent runs per condition (between-conversation stability, small $SD$s, monotonic scaling with intensity). Experiment II: 18-turn dialogues probed at turns 6/12/18 --- self-reports stable, observer-rated expression declines $\sim$20\% with rising variance, attributed to context dilution and RLHF-driven normativity. Persona treated as a validity risk object: drift toward an ``average persona'' would make multi-agent emergence an artifact; proposes reporting standards (mid-session consistency checks, variability metrics, seed/config replication, full documentation) \\
\bottomrule
\end{tabular}
\end{table*}
